\begin{alist}
    \item Έστω ότι το τρίγωνο $\mathrm{A}\mathrm{B}\Gamma$ είναι ισοσκελές με $\mathrm{A}\mathrm{B} = \mathrm{A}\Gamma$ και $\mathrm{B}\mathrm{E}, \Gamma\Delta$ οι διάμεσοι $\mu_\beta, \mu_\gamma$ αντίστοιχα. Θα εξετάσουμε αν είναι $\mu_\beta = \mu_\gamma$.
    
    % TODO: Προσθήκη σχήματος
    
    Τα τρίγωνα $\mathrm{B}\Delta\Gamma$ και $\mathrm{B}\mathrm{E}\Gamma$ έχουν:
    \begin{itemize}
        \item $\mathrm{B}\Delta = \mathrm{E}\Gamma$, ως μισά των ίσων πλευρών $\mathrm{A}\mathrm{B}$ και $\mathrm{A}\Gamma$ του ισοσκελούς τριγώνου $\mathrm{A}\mathrm{B}\Gamma$.
        \item $\mathrm{B}\Gamma$ κοινή πλευρά.
        \item $\hat{\mathrm{B}} = \hat{\Gamma}$, ως γωνίες προσκείμενες στη βάση $\mathrm{B}\Gamma$ του ισοσκελούς τριγώνου $\mathrm{A}\mathrm{B}\Gamma$.
    \end{itemize}
    Σύμφωνα με το Κριτήριο Π-Γ-Π τα τρίγωνα $\mathrm{B}\Delta\Gamma$ και $\mathrm{B}\mathrm{E}\Gamma$ είναι ίσα οπότε θα έχουν και τα υπόλοιπα κύρια στοιχεία τους ίσα, άρα και $\mathrm{B}\mathrm{E} = \Delta\Gamma$, δηλαδή $\mu_\beta = \mu_\gamma$.
    
    \item Η αντίστροφη πρόταση της $\Pi$ διατυπώνεται ως εξής:
    «Αν σε ένα τρίγωνο $\mathrm{A}\mathrm{B}\Gamma$ οι διάμεσοί του $\mu_\beta, \mu_\gamma$ είναι ίσες, τότε το τρίγωνο $\mathrm{A}\mathrm{B}\Gamma$ είναι ισοσκελές με $\beta = \gamma$».
    
    Έστω ότι στο τρίγωνο $\mathrm{A}\mathrm{B}\Gamma$ οι διάμεσοί του $\mu_\beta, \mu_\gamma$, δηλαδή οι $\mathrm{B}\mathrm{E}, \Gamma\Delta$ αντίστοιχα, είναι ίσες. Θα εξετάσουμε αν το τρίγωνο $\mathrm{A}\mathrm{B}\Gamma$ είναι ισοσκελές με $\beta = \gamma$, δηλαδή $\mathrm{A}\mathrm{B} = \mathrm{A}\Gamma$.
    
    % TODO: Προσθήκη σχήματος
    
    Έστω $\mathrm{Z}$ το σημείο τομής των διαμέσων $\mathrm{B}\mathrm{E}$ και $\Gamma\Delta$. Άρα το $\mathrm{Z}$ είναι το βαρύκεντρο του τριγώνου $\mathrm{A}\mathrm{B}\Gamma$ και συνεπώς θα ισχύουν:
    \[ \Delta\mathrm{Z} = \frac{1}{3}\Gamma\Delta \quad (1) \quad \text{και} \quad \mathrm{Z}\mathrm{E} = \frac{1}{3}\mathrm{B}\mathrm{E} \quad (2) \quad \text{αλλά και} \quad \Gamma\mathrm{Z} = \frac{2}{3}\Gamma\Delta \quad (3) \quad \text{και} \quad \mathrm{B}\mathrm{Z} = \frac{2}{3}\mathrm{B}\mathrm{E} \quad (4) \]
    
    Αφού είναι $\Gamma\Delta = \mathrm{B}\mathrm{E}$ (από την υπόθεση), από τις $(1)$ και $(2)$ θα είναι $\Delta\mathrm{Z} = \mathrm{Z}\mathrm{E} \quad (5)$ και από τις σχέσεις $(3)$ και $(4)$ θα είναι $\Gamma\mathrm{Z} = \mathrm{B}\mathrm{Z} \quad (6)$.
    
    Τα τρίγωνα $\Delta\mathrm{Z}\mathrm{B}$ και $\mathrm{E}\mathrm{Z}\Gamma$ έχουν:
    \begin{itemize}
        \item $\Delta\mathrm{Z} = \mathrm{Z}\mathrm{E}$, λόγω της σχέσης $(5)$
        \item $\mathrm{B}\mathrm{Z} = \Gamma\mathrm{Z}$, λόγω της σχέσης $(6)$
        \item $\Delta\hat{\mathrm{Z}}\mathrm{B} = \mathrm{E}\hat{\mathrm{Z}}\Gamma$, ως κατακορυφήν γωνίες
    \end{itemize}
    Επομένως, σύμφωνα με το Κριτήριο Π-Γ-Π τα τρίγωνα $\Delta\mathrm{Z}\mathrm{B}$ και $\mathrm{E}\mathrm{Z}\Gamma$ θα είναι ίσα οπότε θα έχουν και τα υπόλοιπα κύρια στοιχεία τους ίσα. Άρα θα είναι $\Delta\mathrm{B} = \mathrm{E}\Gamma$ οπότε και $2\Delta\mathrm{B} = 2\mathrm{E}\Gamma$ δηλαδή $\mathrm{A}\mathrm{B} = \mathrm{A}\Gamma$.
    Επομένως, το τρίγωνο $\mathrm{A}\mathrm{B}\Gamma$ είναι ισοσκελές.
    
    Συνεπώς, η πρόταση «Αν σε ένα τρίγωνο $\mathrm{A}\mathrm{B}\Gamma$ οι διάμεσοί του $\mu_\beta, \mu_\gamma$ είναι ίσες, τότε το τρίγωνο $\mathrm{A}\mathrm{B}\Gamma$ είναι ισοσκελές με $\beta = \gamma$» είναι αληθής.
    
    \item Εφόσον, από τα ερωτήματα α) και β) η πρόταση $\Pi$ και η αντίστροφή της είναι αληθείς μπορούμε να τις διατυπώσουμε ως την ενιαία πρόταση: «Το τρίγωνο $\mathrm{A}\mathrm{B}\Gamma$ είναι ισοσκελές με $\beta = \gamma$ αν και μόνο αν οι διάμεσοί του $\mu_\beta, \mu_\gamma$ είναι ίσες».
\end{alist}
